% figures/w05-guidance-loop.svg — 바깥 루프가 안쪽 루프에 붙는 자리.
%
%   bash _tools/tikz2svg.sh figures/src/w05-guidance-loop.tex figures/w05-guidance-loop.svg
%
% 사용자 지시 (2026-10-13): "W05에서는 그 다이어그램을 이용해서 outer loop로서
% guidance 명령이 heading 제어기의 desired heading 값이 오고..... 속도 명령도 속도
% 제어 루프에 온다는 것도 연동해서 그려주면 좋을 것 같아"
%
% figures/w03-inner-loop.svg 의 **왼쪽 점선 상자를 연 그림**이다. 그래서 오른쪽
% 절반(안쪽 루프)은 같은 모양·같은 색·같은 블록 이름으로 그리되 **작게** 둔다 —
% 이 그림의 주제는 바깥이고, 안쪽은 3·4 주차에서 이미 열어 보았기 때문이다.
%
% 요점 하나: **바깥 루프의 출력이 곧 안쪽 루프의 명령이다.** 유도는 힘을 만들지
% 않는다. psi_d 를 만들고, 그것을 힘으로 바꾸는 것은 4 주차의 오토파일럿이다.
\documentclass[border=6pt]{standalone}
% gnc-style.tex — 이 볼트의 TikZ 그림이 공유하는 스타일.
%
% 저널 그림 규약을 스타일 하나로 굳혀 둔다. 그림마다 선 굵기나 화살촉을
% 다시 정하지 않는다 — 그것이 그림이 제각각으로 보이는 원인이다.
%
% 쓰는 법:  % gnc-style.tex — 이 볼트의 TikZ 그림이 공유하는 스타일.
%
% 저널 그림 규약을 스타일 하나로 굳혀 둔다. 그림마다 선 굵기나 화살촉을
% 다시 정하지 않는다 — 그것이 그림이 제각각으로 보이는 원인이다.
%
% 쓰는 법:  % gnc-style.tex — 이 볼트의 TikZ 그림이 공유하는 스타일.
%
% 저널 그림 규약을 스타일 하나로 굳혀 둔다. 그림마다 선 굵기나 화살촉을
% 다시 정하지 않는다 — 그것이 그림이 제각각으로 보이는 원인이다.
%
% 쓰는 법:  \input{gnc-style}  (figures/src/ 안에서)

\usepackage{tikz}
\usepackage{amsmath}
\usepackage{xcolor}
\usetikzlibrary{arrows.meta, positioning, calc, fit, backgrounds, decorations.pathmorphing}

% ---- 색 --------------------------------------------------------------------
% 색은 강조 하나에만 쓴다. 나머지는 검정.
\definecolor{ink}{HTML}{1A1A1A}
\definecolor{accent}{HTML}{7C3AED}
\definecolor{warn}{HTML}{B91C1C}
\definecolor{soft}{HTML}{666666}

% 기하 그림(좌표계·유도)에서만 쓰는 넷. 블록선도는 위의 흑백 + accent 로 충분하다.
\definecolor{pathc}{HTML}{0D9488}   % 계획 경로 · 항적
\definecolor{errc}{HTML}{C2410C}    % 오차 (cross-track 등)
\definecolor{alng}{HTML}{B45309}    % along-track
\definecolor{flow}{HTML}{0369A1}    % 조류 · 유속

% ---- 선과 화살촉 -----------------------------------------------------------
\tikzset{
  % 신호선. 화살촉은 작고 뾰족한 것 하나뿐이다.
  sig/.style   = {ink, line width=0.5pt, -{Stealth[length=4pt, width=3pt]}},
  wire/.style  = {ink, line width=0.5pt},
  % 강조 경로 — 파선으로 구분한다. 흑백 인쇄에서도 살아야 하기 때문이다.
  asig/.style  = {accent, line width=0.5pt, densely dashed,
                  -{Stealth[length=4pt, width=3pt]}},
  awire/.style = {accent, line width=0.5pt, densely dashed},
  % 블록. 흰 바탕, 직각, 검은 테두리.
  blk/.style   = {draw=ink, line width=0.55pt, fill=white,
                  minimum width=17mm, minimum height=8mm, inner sep=2pt, font=\small},
  % 합산점
  sum/.style   = {draw=ink, line width=0.55pt, fill=white, circle,
                  inner sep=0pt, minimum size=4.6mm, font=\scriptsize},
  asum/.style  = {draw=accent, line width=0.55pt, fill=white, circle,
                  inner sep=0pt, minimum size=4.6mm, font=\scriptsize},
  % 분기점
  dot/.style   = {ink, fill=ink, circle, inner sep=0pt, minimum size=1.5mm},
  % 신호 이름 — 선 위에 작게
  sname/.style = {font=\small, inner sep=1.5pt},
  % 그림 안의 짧은 설명. 그림 안의 글은 최소로 둔다
  note/.style  = {font=\scriptsize, soft, inner sep=1.5pt},
  anote/.style = {font=\scriptsize, accent, inner sep=1.5pt},
  % 축
  axis/.style  = {ink, line width=0.5pt},
  tick/.style  = {font=\scriptsize, soft},
}

% 캡션 한 줄 — 그림 아래 가는 선과 함께
\newcommand{\gnccaption}[2]{%
  \node[anchor=north west, text width=#1, align=left, font=\scriptsize, ink]
        at (current bounding box.south west) {\vspace{2pt}#2};}
  (figures/src/ 안에서)

\usepackage{tikz}
\usepackage{amsmath}
\usepackage{xcolor}
\usetikzlibrary{arrows.meta, positioning, calc, fit, backgrounds, decorations.pathmorphing}

% ---- 색 --------------------------------------------------------------------
% 색은 강조 하나에만 쓴다. 나머지는 검정.
\definecolor{ink}{HTML}{1A1A1A}
\definecolor{accent}{HTML}{7C3AED}
\definecolor{warn}{HTML}{B91C1C}
\definecolor{soft}{HTML}{666666}

% 기하 그림(좌표계·유도)에서만 쓰는 넷. 블록선도는 위의 흑백 + accent 로 충분하다.
\definecolor{pathc}{HTML}{0D9488}   % 계획 경로 · 항적
\definecolor{errc}{HTML}{C2410C}    % 오차 (cross-track 등)
\definecolor{alng}{HTML}{B45309}    % along-track
\definecolor{flow}{HTML}{0369A1}    % 조류 · 유속

% ---- 선과 화살촉 -----------------------------------------------------------
\tikzset{
  % 신호선. 화살촉은 작고 뾰족한 것 하나뿐이다.
  sig/.style   = {ink, line width=0.5pt, -{Stealth[length=4pt, width=3pt]}},
  wire/.style  = {ink, line width=0.5pt},
  % 강조 경로 — 파선으로 구분한다. 흑백 인쇄에서도 살아야 하기 때문이다.
  asig/.style  = {accent, line width=0.5pt, densely dashed,
                  -{Stealth[length=4pt, width=3pt]}},
  awire/.style = {accent, line width=0.5pt, densely dashed},
  % 블록. 흰 바탕, 직각, 검은 테두리.
  blk/.style   = {draw=ink, line width=0.55pt, fill=white,
                  minimum width=17mm, minimum height=8mm, inner sep=2pt, font=\small},
  % 합산점
  sum/.style   = {draw=ink, line width=0.55pt, fill=white, circle,
                  inner sep=0pt, minimum size=4.6mm, font=\scriptsize},
  asum/.style  = {draw=accent, line width=0.55pt, fill=white, circle,
                  inner sep=0pt, minimum size=4.6mm, font=\scriptsize},
  % 분기점
  dot/.style   = {ink, fill=ink, circle, inner sep=0pt, minimum size=1.5mm},
  % 신호 이름 — 선 위에 작게
  sname/.style = {font=\small, inner sep=1.5pt},
  % 그림 안의 짧은 설명. 그림 안의 글은 최소로 둔다
  note/.style  = {font=\scriptsize, soft, inner sep=1.5pt},
  anote/.style = {font=\scriptsize, accent, inner sep=1.5pt},
  % 축
  axis/.style  = {ink, line width=0.5pt},
  tick/.style  = {font=\scriptsize, soft},
}

% 캡션 한 줄 — 그림 아래 가는 선과 함께
\newcommand{\gnccaption}[2]{%
  \node[anchor=north west, text width=#1, align=left, font=\scriptsize, ink]
        at (current bounding box.south west) {\vspace{2pt}#2};}
  (figures/src/ 안에서)

\usepackage{tikz}
\usepackage{amsmath}
\usepackage{xcolor}
\usetikzlibrary{arrows.meta, positioning, calc, fit, backgrounds, decorations.pathmorphing}

% ---- 색 --------------------------------------------------------------------
% 색은 강조 하나에만 쓴다. 나머지는 검정.
\definecolor{ink}{HTML}{1A1A1A}
\definecolor{accent}{HTML}{7C3AED}
\definecolor{warn}{HTML}{B91C1C}
\definecolor{soft}{HTML}{666666}

% 기하 그림(좌표계·유도)에서만 쓰는 넷. 블록선도는 위의 흑백 + accent 로 충분하다.
\definecolor{pathc}{HTML}{0D9488}   % 계획 경로 · 항적
\definecolor{errc}{HTML}{C2410C}    % 오차 (cross-track 등)
\definecolor{alng}{HTML}{B45309}    % along-track
\definecolor{flow}{HTML}{0369A1}    % 조류 · 유속

% ---- 선과 화살촉 -----------------------------------------------------------
\tikzset{
  % 신호선. 화살촉은 작고 뾰족한 것 하나뿐이다.
  sig/.style   = {ink, line width=0.5pt, -{Stealth[length=4pt, width=3pt]}},
  wire/.style  = {ink, line width=0.5pt},
  % 강조 경로 — 파선으로 구분한다. 흑백 인쇄에서도 살아야 하기 때문이다.
  asig/.style  = {accent, line width=0.5pt, densely dashed,
                  -{Stealth[length=4pt, width=3pt]}},
  awire/.style = {accent, line width=0.5pt, densely dashed},
  % 블록. 흰 바탕, 직각, 검은 테두리.
  blk/.style   = {draw=ink, line width=0.55pt, fill=white,
                  minimum width=17mm, minimum height=8mm, inner sep=2pt, font=\small},
  % 합산점
  sum/.style   = {draw=ink, line width=0.55pt, fill=white, circle,
                  inner sep=0pt, minimum size=4.6mm, font=\scriptsize},
  asum/.style  = {draw=accent, line width=0.55pt, fill=white, circle,
                  inner sep=0pt, minimum size=4.6mm, font=\scriptsize},
  % 분기점
  dot/.style   = {ink, fill=ink, circle, inner sep=0pt, minimum size=1.5mm},
  % 신호 이름 — 선 위에 작게
  sname/.style = {font=\small, inner sep=1.5pt},
  % 그림 안의 짧은 설명. 그림 안의 글은 최소로 둔다
  note/.style  = {font=\scriptsize, soft, inner sep=1.5pt},
  anote/.style = {font=\scriptsize, accent, inner sep=1.5pt},
  % 축
  axis/.style  = {ink, line width=0.5pt},
  tick/.style  = {font=\scriptsize, soft},
}

% 캡션 한 줄 — 그림 아래 가는 선과 함께
\newcommand{\gnccaption}[2]{%
  \node[anchor=north west, text width=#1, align=left, font=\scriptsize, ink]
        at (current bounding box.south west) {\vspace{2pt}#2};}


\definecolor{knownc}{HTML}{1D4ED8}
\definecolor{ctrlc}{HTML}{B45309}

\tikzset{
  blk/.style    = {draw=ink, line width=0.9pt, fill=white, rounded corners=2pt,
                   align=center, inner sep=5pt, font=\small},
  guid/.style   = {blk, draw=accent, line width=1.2pt, fill=accent!5},
  ctrl/.style   = {blk, draw=ctrlc, line width=1.0pt, fill=alng!6, font=\scriptsize},
  alloc/.style  = {blk, draw=knownc, line width=1.0pt, fill=knownc!5, font=\scriptsize},
  plant/.style  = {blk, draw=pathc, line width=1.0pt, fill=pathc!10, font=\scriptsize},
  wp/.style     = {blk, draw=soft, line width=0.8pt, fill=black!3, font=\scriptsize},
  sig/.style    = {ink, line width=0.9pt, -{Stealth[length=5.0pt,width=3.4pt]}},
  cmd/.style    = {accent, line width=1.3pt, -{Stealth[length=5.4pt,width=3.8pt]}},
  fb/.style     = {soft, line width=0.8pt, -{Stealth[length=4.6pt,width=3.2pt]}},
  bus/.style    = {soft, line width=0.8pt},
  lbl/.style    = {font=\scriptsize},
  slx/.style    = {font=\scriptsize\ttfamily, soft},
  capt/.style   = {font=\scriptsize\bfseries, warn},
  band/.style   = {font=\small\bfseries, ink},
}

\begin{document}
\begin{tikzpicture}[x=1cm, y=1cm]

% =====================================================================
% 바깥 루프 — 웨이포인트에서 두 명령까지
% =====================================================================
\node[wp, text width=1.9cm] (WP) at (-0.2,0.55)
  {\textbf{WAYPOINTS}\\[2pt] $(N_k, E_k)$\\ $k = 1 \ldots n$};

\node[guid, text width=4.6cm] (GD) at (4.1,0.55)
  {\textbf{GUIDANCE} \ \ \emph{outer loop}\\[4pt]
   {\scriptsize 1 \ \ leg direction \ $\pi_p = \mathrm{atan2}(\Delta E, \Delta N)$}\\[2pt]
   {\scriptsize 2 \ \ rotate into the leg \ $(x_e,\, y_e)$}\\[2pt]
   {\scriptsize 3 \ \ switch when $d - x_e < R$}\\[2pt]
   {\scriptsize 4 \ \ the law \ $\psi_d = \pi_p - \arctan\!\dfrac{y_e + \kappa y_{\mathrm{int}}}{\Delta}$}};
\node[slx, anchor=north, yshift=-1pt] at (GD.south) {guidance};
\draw[sig] (WP.east) -- (GD.west);

\node[wp, text width=1.9cm] (UD) at (-0.2,-2.6)
  {\textbf{SPEED ORDER}\\[2pt] $u_d$ \ \ chosen by\\ the mission};

% =====================================================================
% 두 명령 — 이 그림의 요점
% =====================================================================
\draw[cmd] (GD.east) -- (8.5,0.55) -- (8.5,0.05) -- (9.15,0.05);
\node[lbl, accent, anchor=south] at (7.5,0.62) {$\psi_d$};
\draw[cmd] (UD.east) -- (8.5,-2.6) -- (8.5,-1.95) -- (9.15,-1.95);
\node[lbl, accent, anchor=south] at (8.0,-2.52) {$u_d$};

% =====================================================================
% 안쪽 루프 — 3·4 주차가 이미 연 상자들. 작게 둔다
% =====================================================================
\node[ctrl, text width=2.5cm] (HD) at (10.6,0.05)
  {\textbf{HEADING LOOP}\\[2pt] $N = K_p\,\mathrm{ssa}(\psi_d-\psi) - K_d r$};
\node[slx, anchor=north, yshift=-1pt] at (HD.south) {heading autopilot};

\node[ctrl, text width=2.5cm] (SP) at (10.6,-1.95)
  {\textbf{SPEED LOOP}\\[2pt] $X = K_{p,u}(u_d-u) + I$};
\node[slx, anchor=north, yshift=-1pt] at (SP.south) {speed loop};

\node[alloc, text width=2.4cm] (AL) at (14.3,-0.95)
  {\textbf{ALLOCATION}\\[3pt]
   $\begin{bmatrix} F_L \\ F_R \end{bmatrix} = \mathbf{T}^{-1}\begin{bmatrix} X \\ N \end{bmatrix}$\\[3pt]
   $n = \pm\sqrt{\lvert F\rvert/k}$};
\node[slx, anchor=north, yshift=-1pt] at (AL.south) {allocation};

\node[plant, text width=1.5cm] (OT) at (17.9,-0.95)
  {\textbf{OTTER}\\[1pt] equations\\ of motion};
\node[slx, anchor=north, yshift=-1pt] at (OT.south) {Otter};

\draw[sig] (HD.east) -- (12.9,0.05) -- (12.9,-0.55) -- (AL.west |- 0,-0.55);
\node[lbl, anchor=south] at (12.35,0.12) {$N$};
\draw[sig] (SP.east) -- (12.9,-1.95) -- (12.9,-1.35) -- (AL.west |- 0,-1.35);
\node[lbl, anchor=north] at (12.35,-2.02) {$X$};
\draw[sig] (AL.east) -- (OT.west);
\node[lbl, anchor=south] at (16.3,-0.88) {$n_L,n_R$};
\draw[sig] (OT.east) -- (19.4,-0.95);
\node[lbl, anchor=west] at (19.45,-0.95) {$\mathbf{x}$};

% =====================================================================
% 되먹임 둘 — 안쪽은 짧고 빠르게, 바깥은 길고 느리게
% =====================================================================
\draw[bus] (19.75,-0.95) -- (20.2,-0.95) -- (20.2,-4.35) -- (9.6,-4.35);
\fill[soft] (19.25,-0.95) circle (1.4pt);
\draw[fb] (9.6,-4.35) -- (9.6,-2.52);
\node[lbl, soft, anchor=west] at (9.67,-3.55) {$u$};
\draw[fb] (12.0,-4.35) -- (12.0,-0.52);
\fill[soft] (12.0,-4.35) circle (1.4pt);
\node[lbl, soft, anchor=west] at (12.07,-3.55) {$\psi,\ r$};
\node[capt, anchor=west] at (13.4,-4.20) {INNER feedback --- fast};

\draw[bus] (20.2,-4.35) -- (20.2,-5.85) -- (2.0,-5.85);
\draw[fb] (2.0,-5.85) -- (2.0,1.65) -- (GD.west |- 0,1.65);
\node[lbl, soft, anchor=south] at (6.0,-5.80) {position $(N, E)$ and heading $\psi$};
\node[capt, anchor=west] at (13.4,-5.72) {OUTER feedback --- slow};

% =====================================================================
% 띠 설명
% =====================================================================
\node[band, anchor=west] at (-1.4,2.75)
  {OUTER LOOP \ --- \ guidance makes commands, never forces};
\draw[bus] (-1.3,2.45) -- (8.3,2.45);
\node[band, anchor=west, ctrlc] at (9.2,2.75) {INNER LOOP \ --- \ Weeks 3 and 4};
\draw[bus] (9.3,2.45) -- (18.6,2.45);

%  이 문장이 그림의 요점이다. 띠 사이에 두면 INNER 띠와 겹친다 — 유도 상자 아래로.
\node[capt, anchor=north, align=center] at (4.1,-1.45)
  {the output of the outer loop \emph{is} the command of the inner loop};

\end{tikzpicture}
\end{document}
